\documentclass[12pt]{article}

\usepackage[utf8]{inputenc}
\usepackage{latexsym,amsfonts,amssymb,amsthm,amsmath}
\usepackage{tikz}
\usepackage{stanli}
\usetikzlibrary{shapes,calc,through,3d,patterns}
\usetikzlibrary{arrows.meta,decorations.pathreplacing,angles,quotes}
\usepackage{multicol}
\usepackage{geometry}
\usepackage{mathtools}
\usepackage{siunitx}
\usepackage{booktabs}

\setlength{\parindent}{0in}
\setlength{\oddsidemargin}{-0.25in}
\setlength{\textwidth}{7in}
\setlength{\textheight}{9.3in}
\setlength{\topmargin}{-1in}
\setlength{\headheight}{18pt}

\newcommand\blfootnote[1]{%
    \begingroup
    \renewcommand\thefootnote{}\footnote{#1}%
    \addtocounter{footnote}{-1}%
    \endgroup
}

\definecolor{darkred}{rgb}{0.6,0,0}

\newcommand{\bolddelta}{\boldsymbol\delta}
\newcommand{\boldxi}{\boldsymbol\xi}

\begin{document}

\noindent ME473 - Dynamic finite element analysis of structures\hfill EPFL\\
Stefano Burzio \hfill 2025\\
\vspace{0.3cm}

\hspace*{\fill}\textbf{\Large Problem set 9}\hspace*{\fill}

\hrulefill

\subsection*{Problem 1}
Use the Rayleigh minimization principle to determine the natural frequencies of a free-free shaft in torsion, assuming the structure is discretized using a single linear finite element. The strong form is given by
$$ \frac{\partial}{\partial x}\left(GI_p \frac{\partial}{\partial x}\varphi(x, t)\right) = \rho I_p \ddot{\varphi}(x, t)$$
where the shaft parameters are:
\begin{multicols}{2}
  \begin{itemize}
    \item length $\ell$,
    \item polar moment of inertia $I_p$,
    \item shear modulus $G$,
    \item mass density $\rho$.
  \end{itemize}
\end{multicols}
The unknown $\varphi$ represent the torsional vibration: angular displacement (twist angle) at position $x \in [0, \ell]$ and time $t \in [0, T]$. At $x = 0$ and $x = \ell$ there are no external applied moments, this translates into natural boundary conditions:
$$
  \frac{\partial}{\partial x}\varphi(0,t) = \frac{\partial}{\partial x}\varphi(\ell,t) = 0 \qquad \forall t\in]0,T[.
$$

\subsection*{Problem 2}
Develop a \texttt{MATLAB} script to compute the first two natural frequencies and the corresponding mode shapes associated with the longitudinal vibrations of a uniform free-free bar, using the subspace iteration method. The structure is characterized by the following physical parameters:
\begin{multicols}{2}
  \begin{itemize}
    \item length: $\ell=0.5$ m,
    \item cross-sectional area: $A=4\,10^{-4}$ m$^2$,
    \item Young's modulus: $E=210$ GPa,
    \item mass density: $\rho=7850$ kg/m$^3$.
  \end{itemize}
\end{multicols}
The bar is modeled using a single quadratic finite element with three equally spaced nodes. The structure stiffness and lumped mass\footnote{The use of a lumped mass matrix is often justified by its ability to simplify computations, improve numerical stability in explicit time integration schemes, and avoid non-physical oscillations associated with consistent mass formulations in low-order elements.} matrices are given by:
\[
  \mathbf{K} = \frac{EA}{3\ell}
  \begin{bmatrix}
    7  & -8 & 1  \\
    -8 & 16 & -8 \\
    1  & -8 & 7
  \end{bmatrix}\qquad\text{and}\qquad
  \mathbf{M} = \frac{\rho A \ell}{6} \text{diag}(1, 4, 1).
\]

\begin{enumerate}
  \item[1)] Formulate the generalized eigenvalue problem for the system, and define the stiffness and mass matrices explicitly.

  \item[2)] Since the system exhibits a rigid-body mode with zero frequency, a spectral shift is applied to ensure convergence. The modified eigenvalue problem is:
        \[
          (\mathbf{K} + \sigma \mathbf{M}) \mathbf{p} = (\lambda + \sigma) \mathbf{M} \mathbf{p}
        \]
        where the shift parameter is chosen (to simplify the computation) as:
        \[
          \sigma = 2\frac{E}{\rho \ell^2}.
        \]
        Define the shifted stiffness matrix as $\mathbf{K}_\sigma = \mathbf{K} + \sigma \mathbf{M}$, yielding the equivalent eigenvalue problem:
        \[
          \mathbf{K}_\sigma \mathbf{p} = \lambda_\sigma \mathbf{M} \mathbf{p}, \quad \text{with} \quad \lambda_\sigma = \lambda + \sigma.
        \]

  \item[3)] Implement three iterations of the subspace iteration algorithm to approximate the two lowest eigenpairs $(\lambda_i, \mathbf{p}_i)$. For simplicity use the function \texttt{eig} to solve the reduced Rayleigh's minimization problem. Use the following initial guess:
        \[
          \mathbf{P}_{(0)} = \frac{1}{\sqrt{\rho A \ell}} \begin{bmatrix}
            1 & 0 \\
            0 & 1 \\
            0 & 0
          \end{bmatrix}.
        \]

  \item[4)] Compare the numerical approximation with the exact (shifted) eigenvalues and corresponding eigenvectors (defined up to a multiplicative constant):
        \[
          \lambda_{\sigma,1} = 2\frac{E}{\rho \ell^2}\qquad\text{and}\qquad \lambda_{\sigma,2} = 14\frac{E}{\rho \ell^2},
        \]
        \[
          \mathbf{p}_1 = \begin{bmatrix} -1 \\ -1 \\ -1 \end{bmatrix}\qquad\text{and}\qquad
          \mathbf{p}_2 = \begin{bmatrix} -1 \\ 0 \\ 1 \end{bmatrix}
        \]
        and evaluate the accuracy by computing:
        \begin{itemize}
          \item The relative error between approximated and exact eigenvalues.
          \item The error in the $\mathbf{M}$-norm between approximated and exact eigenvectors.
        \end{itemize}

        %\item Investigate the sensitivity of convergence to the choice of initial vectors. Repeat the computation using the alternative initial matrix:
        %   \[
        %    \mathbf{P}_{(0)} = \frac{1}{\sqrt{\rho A \ell}} \begin{bmatrix}
        %      1 & 0 \\
        %      4 & 1 \\
        %      1 & 0
        %    \end{bmatrix}
        %  \]
        %  and discuss the observed changes in convergence behavior.
\end{enumerate}
\end{document}